\documentclass[11pt,a4paper]{article}
\usepackage[T1]{fontenc}
\usepackage[utf8]{inputenc}
\usepackage[french]{babel}
\usepackage{lmodern}
\usepackage[a4paper,margin=2.2cm]{geometry}
\usepackage{booktabs,tabularx}
\usepackage[hidelinks]{hyperref}
\usepackage{fancyhdr}
\setlength{\parindent}{0pt}
\setlength{\parskip}{0.55em}
\pagestyle{fancy}
\fancyhf{}
\lhead{BD70 - Vélo Libre-Service}
\rhead{Rapport de performance}
\cfoot{\thepage}
\begin{document}

\begin{center}
{\LARGE\textbf{Rapport d'analyse des performances}}\\[0.4em]
{\large Système de vélos en libre-service}\\[0.8em]
Floryan Bornet \textperiodcentered\ Tim Carrara \textperiodcentered\ Enzo Vidy\\
18 novembre 2026
\end{center}

\section{Objectif}
Ce rapport mesure l'effet d'un index composé sur une requête centrale de l'application : l'affichage de l'historique de position d'un vélo. Cette consultation filtre un vélo donné puis ordonne ses déplacements du plus récent au plus ancien. Elle est utilisée par la vue opérateur ``Historique des vélos''.

\section{Jeu de données et protocole}
Les mesures ont été réalisées sur MySQL local, le 18 novembre 2026, dans une base temporaire recréée à partir des scripts livrés. Le jeu de charge ajoute 150 000 positions à \texttt{VELO-001}. La base mesurée contenait donc 153 015 lignes dans \texttt{historique\_position}, en plus de 3 002 réservations, 3 017 changements d'état, 12 usagers, 12 vélos et 32 bornes.

Le script \texttt{benchmark\_optimisation.sql} mesure la même requête dans les deux scénarios. Au-delà de 50 000 lignes, il réalise trois répétitions et conserve les 500 dernières positions retournées à chaque itération. Le temps est calculé par \texttt{TIMESTAMPDIFF(MICROSECOND, ...)} puis exprimé en millisecondes.

\begin{verbatim}
SELECT h.id_historique_position
FROM historique_position h
LEFT JOIN borne b_depart ON b_depart.id_borne = h.ancienne_borne
LEFT JOIN borne b_arrivee ON b_arrivee.id_borne = h.nouvelle_borne
WHERE h.numero_serie = 'VELO-001'
ORDER BY h.date_heure DESC
LIMIT 500;
\end{verbatim}

\section{Optimisation étudiée}
Le scénario initial conserve un index simple sur \texttt{numero\_serie}. Il est utile pour filtrer le vélo mais ne correspond pas entièrement au besoin de tri chronologique. Le scénario optimisé utilise l'index composé suivant :
\begin{verbatim}
CREATE INDEX idx_historique_position_velo_date
ON historique_position (numero_serie, date_heure);
\end{verbatim}

L'ordre des colonnes est volontaire : la première correspond au prédicat d'égalité \texttt{WHERE numero\_serie = ...}, puis la seconde à \texttt{ORDER BY date\_heure DESC}. Les enregistrements du même vélo sont donc structurés par date dans l'index.

\section{Résultats}
\begin{center}
\begin{tabular}{lrrrr}
\toprule
\textbf{Scénario} & \textbf{Répétitions} & \textbf{Temps total} & \textbf{Moyenne} & \textbf{Gain}\
\midrule
Index simple sur \texttt{numero\_serie} & 3 & 799,076 ms & 266,359 ms & référence\\
Index composé \texttt{(numero\_serie, date\_heure)} & 3 & 22,176 ms & 7,392 ms & $\times$36,0\\
\bottomrule
\end{tabular}
\end{center}

Le gain observé est d'environ 258,967 ms par exécution, soit une durée moyenne divisée par 36. Le résultat est cohérent avec l'objectif de l'index : réduire le travail nécessaire après le filtrage du vélo et servir directement un ordre compatible avec la chronologie demandée.

\section{Vérification du plan}
Le \texttt{EXPLAIN FORMAT=JSON} produit par le script indique que MySQL sélectionne \texttt{idx\_historique\_position\_velo\_date} dans le scénario optimisé. Les jointures vers les bornes et stations s'appuient sur leurs clés primaires avec le type d'accès \texttt{eq\_ref}. Le point coûteux est donc bien l'accès à l'historique, qui est précisément celui visé par l'index.

\section{Autres index ciblés}
La base ne crée pas un index sur chaque colonne : un index accélère les lectures mais consomme de l'espace et doit être maintenu lors des insertions et mises à jour. Les index ont été choisis à partir des requêtes fréquentes :
\begin{itemize}
  \item \texttt{reservation(statut, date\_reservation)} pour l'expiration périodique des réservations ;
  \item \texttt{trajet(id\_usager, date\_depart)} pour l'historique d'un usager ;
  \item \texttt{trajet(numero\_serie, date\_arrivee)} pour vérifier qu'un vélo n'est pas déjà en location ;
  \item \texttt{borne(id\_station, etat)} et \texttt{velo(etat, id\_categorie)} pour l'état des stations ;
  \item \texttt{historique\_etat(numero\_serie, date\_heure)} pour la chronologie des changements d'état.
\end{itemize}

\section{Limites et reproductibilité}
Un temps absolu dépend de la machine, de la version de MySQL, de la mémoire disponible et du cache. La conclusion fiable est donc la comparaison sur la même machine, avec la même requête et le même jeu de données. Pour reproduire cette mesure, importer les scripts de schéma et de données, charger \texttt{donnees\_charge\_optimisation.sql}, puis exécuter \texttt{benchmark\_optimisation.sql}. Les temps affichés par la table temporaire \texttt{resultat\_benchmark} sont ceux à présenter pour l'environnement utilisé.

\end{document}
